\chapter{Model description}
During the years the effectiveness of graph-based models in fraud detection brought them to be the standard approach in the field \cite{artgrafi}.
For our project we will consider a financial transaction graph $G = (V, E)$, \(|V| \doteq N\), where each node $v \in V$ represents an account and each edge $(u, v) \in E$ represents a transactions between accounts $u$ and $v$. Each edge is associated with a weight $W(u,v) = w_{uv}$ defined as
\begin{equation}
    w_{uv} \doteq \sum_{t \in T_{uv}} \operatorname{amount}(t)
\end{equation}
where $T_{uv}$ is the set of transactions between accounts $u$ and $v$. Node centrality is also taken in consideration and it's computed as
\begin{equation}
    \alpha_{i} \doteq \frac{\operatorname{deg}(i)}{\sum_{j\in V}\operatorname{deg}(j)}
\end{equation}
Before moving forward the graph is filtered to remove edges with low weights
\begin{equation}
    E \leftarrow \{(u,v) \in E \; | \; w_{uv} > \tau_\text{filter}\}
\end{equation}
This is true for every dataset mentioned in Section \ref{sc:objectives} but for the Bitcoin transaction graph where nodes are transactions and edges are the flow of Bitcoin.

\section{Generalized Amplitude Encoding}
Generalized Amplitude Encoding with Entanglement Seeding (GAES) is performed to encode the graph structure and features into quantum states.
The idea behind standard amplitude encoding is to represent data as amplitudes of a quantum state. Given a vector \(v \in \mathbb{R}^{N}, N=2^n\) with \(\|v\| = 1\), the corresponding quantum state is
\begin{equation}
    |\psi\rangle = \sum_{i=0}^{2^n-1} v_i |i\rangle
\end{equation}
where \(\psi\) is a quantum state with \(n\) qubits.

In our case we want to encode not only node features but also graph structure. We achive this by usingi entanglement seeding \cite{mainarticle}, which introduces entanglement between qubits to capture relationships between nodes.
\begin{equation}\label{statedef}
    |\psi_{G}\rangle=\frac{1}{\sqrt{Z}}\sum_{i,j}A_{ij}\mathcal{E}_{ij}(\theta_{e})|i\rangle|j\rangle
\end{equation}
where
\begin{equation}
    A_{ij} \doteq \begin{cases}
        w_{ij}   & \text{if } i \neq j \\
        \alpha_i & \text{if } i = j
    \end{cases}
\end{equation}
\begin{equation}
    Z \doteq \sum_{i,j} |A_{ij}|^{2}
\end{equation}
and
\begin{equation}
    \mathcal{E}_{ij}(\theta_{e}) \doteq \exp(-i\theta_{e}\sigma_{i}^{x}\sigma_{j}^{x})
\end{equation}
is a quantum gate that applies entanglement between qubits $i$ and $j$, \(\sigma_i^A\) with \(A \in \{X, Y, Z\}\) are the Pauli operators acting on qubit \(i\) and $\theta_{e}$ are parameters that will be optimized during training so that the state physically embodies the interactions defined by $H(G)$.

This definition is, however, subject to a big dimensionality flaw. Under our amplitude encoding, the quantum state lives in an $2n$-qubit Hilbert space where $n = \log_2(N)$. The operator $\sigma_i^x$ is a global Pauli-X gate acting on physical qubit $i$, yet here the index $i$ ranges over node indices $[0, N-1]$, which is dimensionally inconsistent with the $n$-qubit register. More critically, even if the index were in range, applying $\sigma^x$ based on a binary mapping would universally flip bits across the entire superposition, completely scrambling the graph topology instead of selectively linking node $i$ to node $j$.

My solution is to replace the entanglement operator with a \textit{subspace projector} (exchange operator) that acts exclusively on the subspace spanned by the computational basis states $|i\rangle$ and $|j\rangle$:
\begin{equation}
    \mathcal{E}_{ij}(\theta_{e}) \doteq \exp\bigl(-i\theta_{e}(|i\rangle\langle j| + |j\rangle\langle i|)\bigr)
\end{equation}
The generator $(|i\rangle\langle j| + |j\rangle\langle i|)$ is the generator of a continuous-time quantum walk between nodes $i$ and $j$. When exponentiated, it applies an $SU(2)$ rotation strictly to the amplitudes of those two connected nodes, leaving all other amplitudes untouched. This is consistent with the amplitude encoding paradigm and achieves the authors' intended goal of edge-specific entanglement.

\(H(G)\) is the graph Hamiltonian and it includes higher-order interactions to capture multi-party transactions. The original paper define it as
\begin{equation}\label{H}
    H(G) \doteq \sum_{(i,j)\in E}w_{ij}(\sigma_{i}^{x}\sigma_{j}^{x}+\sigma_{i}^{y}\sigma_{j}^{y}+\sigma_{i}^{z}\sigma_{j}^{z})+\sum_{i\in V}\alpha_{i}\sigma_{i}^{z}+\sum_{(i,j,k)\in\Delta}\gamma_{ijk}\sigma_{i}^{z}\sigma_{j}^{z}\sigma_{k}^{z}
\end{equation}
where $\gamma_{ijk}$ captures the three-way interactions between nodes \(i\), \(j\), and \(k\). The original paper doesn't provide a specific definition for \(\gamma_{ijk}\), just the description of its purpose. A possible interpretation is the geometric mean of the weights of the edges forming the triangle.
\begin{equation}
    \gamma_{ijk} \doteq \sqrt[3]{w_{ij} \cdot w_{jk} \cdot w_{ki}}
\end{equation}
From my prespective there is a big flaw in the definition of the Hamiltonian. As before, from equation \ref{H} we can see that the Paoli gates are being applied to the \(i\)-th qubits, this index \(i\) layes in the range \([0, N-1]\) where \(N\) is the number of nodes in the graph. But from equation \ref{statedef} we can see that after the encoding we will have \(2n\) qubits, where \(n = \log_2(N)\). This is a contradiction and obviously we cannot afford to have \(N\) qubits since in real world applications \(N\) can be in the order of millions.

My solution is to redefine the Hamiltonian using lernable parameters that will capture the interactions between nodes without explicitly mapping them to qubits \cite{newHam}. The new Hamiltonian becomes
\begin{equation}
    H \doteq H_{eff}(\Theta) = \underbrace{\sum_{q=0}^{n-1} \phi_q \sigma_q^z}_{\text{Global Bias}} + \underbrace{\sum_{q < r}^{n-1} \theta_{qr} \sigma_q^z \sigma_r^z}_{\text{Qubit Correlations}} + \underbrace{\sum_{q < r < s}^{n-1} \psi_{qrs} \sigma_q^z \sigma_r^z \sigma_s^z}_{\text{Higher-Order Entanglement}}
\end{equation}
This consists in a parameterized diagonal Hamiltonian that strictly operates on the correct $2n$-qubit register, making the Variational Quantum Circuit (VQC) mathematically executable. This happens at the expense of the loss of explicitly topology. In the original, albeit flawed concept, the Hamiltonian explicitly used edge weights $w_{ij}$ to drive the quantum interactions. With this new implementation we are assuming that the information contained in the quantum stare \(\psi_G\) is sufficient to reconstruct the graph properties and that the parameters \(\Theta\) will learn the topological information from the data.

After the encoding we will end up with a quantum state \(\psi_G\) composed by \(2n\) qubits, where \(n \doteq \log_2(N)\). Usually for amplitude encoding we would need only \(n\) qubits, the other \(n\) are defined as \textit{enviorment} and are introduced to perform the entaglement seeding. In the equation \ref{statedef} those correspond to the second register \(|j\rangle\).
\begin{equation}
    |\psi_G\rangle \in \mathcal{H}_\text{source} \otimes \mathcal{H}_\text{env}
\end{equation}
Since we are interested only in the source register, we will trace out the environment qubits to obtain a reduced density matrix \(\rho_G\). The use of the density matrix alows to represent mixed states and capture statistical properties of the quantum system.
\begin{equation}
    \rho_\text{total} = |\psi_G\rangle\langle\psi_G|
\end{equation}
\begin{equation}
    \rho_G = \operatorname{Tr}_{\text{env}}(\rho_\text{total})
\end{equation}
where
\begin{equation}
    \operatorname{Tr}_{\text{env}}(\rho_\text{total}) \doteq \sum_{m \in \text{env}} (\mathbb{I} \otimes \langle m|) \rho_{\text{total}} (\mathbb{I} \otimes |m\rangle)
\end{equation}
In the end we will obtain a quantum state \(\rho_G\) composed by \(n\) qubits that encapsulates the information of the original graph \(G\) by effectively incorporating the entanglement effects introduced during the encoding process. Although the final representation needs only \(n\) qubits, a quantum computer with \(2n\) qubits will be needed to perform the encoding.

\section{Variational Quantum Graph Convolution}
During this step parameterized quantum circuits perform non-linear graph convolutions, incorporating higher-order neighborhood interactions to model transaction flows evolving through the layers with the non linear operator
\begin{equation}
    \rho^{(0)}=\rho_{G}
\end{equation}
\begin{equation}
    \rho^{(l+1)}=\mathcal{N}_{\theta}^{(l)}(U_{\theta}^{(l)}\rho^{(l)}U_{\theta}^{(l)\dagger})
\end{equation}
where \(\mathcal{N}_\theta^{(l)}(\rho)\) is the non-linear quantum channel defined as
\begin{equation}
    \mathcal{N}_{\theta}^{(l)}(\rho)=\sum_{k}p_{k}(\theta)\Pi_{k}\rho\Pi_{k}
\end{equation}
\(U_{\theta}^{(l)}\) is the unitary operator. The original paper defines it as
\begin{equation}
    U_{\theta}^{(l)}=\operatorname{exp}\left(-i\sum_{(i,j)\in E}\theta_{ij}^{(l)}H_{ij}-i\sum_{i\in V}\phi_{i}^{(l)}\sigma_{i}^{z}-i\sum_{(i,j,k)\in\Delta}\psi_{ijk}^{(l)}\sigma_{i}^{z}\sigma_{j}^{z}\sigma_{k}^{z}\right)
\end{equation}

This definition is subject to the same dimensionality flaw already identified for the Hamiltonian $H(G)$. Summing $\sigma_i^z$ over all vertices $V$ and $\sigma_i^z\sigma_j^z\sigma_k^z$ over all hyperedges $\Delta$ implicitly assumes $N$ distinct qubits, which is incompatible with the $n$-qubit amplitude-encoding space.

Following the same reasoning used for the Hamiltonian, my solution is to redefine the variational unitary using the effective Hamiltonian $H_\text{eff}$ \cite{artmodificheencoding}, now parameterized per layer:
\begin{equation}
    U_{\theta}^{(l)} = \exp\bigl(-i\, H_\text{eff}(\Theta^{(l)})\bigr)
\end{equation}
where the per-layer effective Hamiltonian is
\begin{equation}
    H_\text{eff}(\Theta^{(l)}) = \sum_{q=0}^{n-1} \phi_q^{(l)} \sigma_q^z + \sum_{q<r}^{n-1} \theta_{qr}^{(l)} \sigma_q^z \sigma_r^z + \sum_{q<r<s}^{n-1} \psi_{qrs}^{(l)} \sigma_q^z \sigma_r^z \sigma_s^z
\end{equation}
This substitution makes the VQC mathematically executable on NISQ hardware. However, as with the Hamiltonian correction, it comes at a cost: the model no longer performs explicit topological message passing in the hidden layers. The VQC effectively acts as a non-linear phase classifier, and the structural information of the graph must be fully carried by the initial state $\rho_G$ produced during the GAES encoding step.

In the original paper the Kraus operators \(\Pi_k\) and the probabilities \(p_k(\theta)\) are not explicitly defined. In my implementation, in regard to preserve the node information (diagonal terms of the density matrix) while filtering the complex relational patterns, i chose to implement the \textit{Phase-Damping (Dephasing) Channel}. This should break unitarity, enabling the model to learn complex, non-linear boundaries between legitimate and fraudulent transactions \cite{phasedamp}.
The mathematical definition of Kraus operators becomes
\begin{equation}
    \Pi_1 = I = \begin{pmatrix} 1 & 0 \\ 0 & 1 \end{pmatrix}, \quad \Pi_2 = \sigma^z = \begin{pmatrix} 1 & 0 \\ 0 & -1 \end{pmatrix}
\end{equation}
Notice that they both satisfy the completeness relation
\begin{equation}
    \sum_{k} E_k^\dagger E_k = I
\end{equation}
The associated probabilities depend on the tunable parameters \(\theta_1, \theta_2\) and are defined as
\begin{equation}
    p_0(\theta) = \frac{e^{\theta_0}}{e^{\theta_0} + e^{\theta_1}}, \quad p_1(\theta) = \frac{e^{\theta_1}}{e^{\theta_0} + e^{\theta_1}}
\end{equation}
The bigger will be the value of \(\theta_1\) the more noise will be added to the system. \(\theta\) should be initialized in a way that \(p_0(\theta) \gg p_1(\theta)\), the network will learn to increase the "noise" (phase damping) in deeper layers to strip away irrelevant correlations and focus on the core topological features of the fraud graph.

$\theta_{ij}^{(l)}$, $\phi_{i}^{(l)}$ and $\psi_{ijk}^{(l)}$ are tunable parameters that captures transaction dynamics, account-specific dynamics and multi-party (hyperedge) interactions respectively. The quantum feature embedding after $L$ layers will be
\begin{equation}
    Z_{Q}=Tr_{anc}\left(\prod_{l=1}^{L}\mathcal{N}_{\theta}^{(l)}(U_{\theta}^{(l)}\rho_{G}U_{\theta}^{(l)\dagger})\right)
\end{equation}
Where \(Tr_{\text{anc}}\) is defined as for \(Tr_{\text{env}}\) and it traces out ancilla qubits that were introduced to apply the Kraus operators.

We also compute the correlation entropy which captures multi-scale transaction interactions
\begin{equation}
    C_{Q}=\sum_{i\ne j}Tr(\rho_{ij}^{(L)}log~\rho_{ij}^{(L)}-\rho_{i}^{(L)}log~\rho_{i}^{(L)}-\rho_{j}^{(L)}log~\rho_{j}^{(L)})
\end{equation}
where $\text{Tr}(\rho) = \sum_k \rho_{kk}$.

Both \(Z_{Q}\) and \(C_{Q}\) will be used as features in the final classification.

\section{Topological Quantum Signature Extraction}
In this section we will extract Quantum Betti numbers and persistent Euler characteristic from the quantum-encoded graph.
\begin{definition}[Betti numbers]
    Betti numbers act as a "counter" for different types of holes in a topological space or object. Each Betti number \(\beta_i\) corresponds to the number of i-cells in the CW-complex (graph). In particular, \(\beta_0\) is the number of connected components, \(\beta_1\) counts the number of one-dimensional holes (Cycles), \(\beta_2\) counts the number of two-dimensional voids (cavities). In our methods higher-order Quantum Betti numbers are computed to extend the analysis beyond simple loops.
\end{definition}
\begin{definition}[Euler's characteristic]
    Let \(\tau\) be a CW-complex composed by \(\beta_i\) i-cells. We define Euler's characteristic as
    \begin{equation}
        \chi(\tau) = \sum_{i=0}^{\infty} (-1)^i \beta_i
    \end{equation}
    The Euler's characteristic is a topological invariant that combines information from all Betti numbers into a single scalar value.
\end{definition}
The original paper doesn't specify how many Betti numbers are computed, in my implementation I decided to compute up to \(\beta_3\) to capture complex topological features.

From now on we define \(\rho \doteq \rho^{(L)}\) from the last layer of the Variational Quantum Graph Convolution. Persistent homology is extended to compute higher-order topological invariants over a multi-scale quantum distance landscape, capturing structural anomalies in transaction networks. The quantum distance matrix uses a weighted metric defined as follows.
\begin{equation}
    d_{\rho}(i,j)=\sqrt{1-F_{w}(\rho_{i},\rho_{j})}
\end{equation}
where
\begin{equation}
    F_{w}(\rho_{i},\rho_{j})=\frac{(Tr\sqrt{\sqrt{\rho_{i}}W\rho_{j}\sqrt{\rho_{i}}})^{2}}{Tr(W)}
\end{equation}
and \(W\) is a positive definite weight matrix encoding priors like transaction frequency and account centrality.
The quantum Vietoris--Rips complex is defined as the simplicial complex formed by connecting nodes that are within a distance $\epsilon$ of each other.
\begin{equation}
    \mathcal{VR}_{\epsilon}(\rho)=\{\sigma\subseteq V \mid d_{\rho}(i,j)\le\epsilon,\forall i,j\in\sigma\}
\end{equation}
Then the Quantum Betti numbers \(\beta_k(\epsilon)\) are computed as
\begin{equation}
    \beta_{k}^{Q}(\epsilon)=\dim H_{k}(\mathcal{V}\mathcal{R}_{\epsilon}(\rho);\mathbb{Q})
\end{equation}
where \(H_k\) is the k-th homology group with rational coefficients defined as
\begin{equation}
    H_{k}(\mathcal{VR}_{\epsilon}(\rho);\mathbb{Q})=\frac{\ker \partial_{k}}{\operatorname{im} \partial_{k+1}}
\end{equation}
with \(\partial_k : C_k \to C_{k-1}\) is the boundary operator, a linear map that takes a $k$-dimensional shape and returns its $(k-1)$-dimensional boundary.

Once the Quantum Betti numbers are computed we can define the Persistent Euler Characteristic (PEC) as
\begin{equation}
    \chi^{Q}(\epsilon)=\sum_{k}(-1)^{k}\beta_{k}^{Q}(\epsilon)
\end{equation}
The value of \(\epsilon\) should be tuned and, since \(d_\rho(i,j) \in [0,\sqrt{2}]\), we should choose values in that range. More then one value can be choosed and used as features for the final classification.
Adding those topological features to the quantum features \(Z_Q\) and \(C_Q\) we obtain the final feature vector
\begin{equation}
    \phi(G)=[Z_{Q},C_{Q},\{\beta_{k}^{Q}(\epsilon)\}_{k,\epsilon},\{\chi^{Q}(\epsilon)\}_{\epsilon}]
\end{equation}
\section{Quantum-Classical Hybrid Anomaly Learning}
In this section a hybrid quantum-classical architecture is proposed to perform the final anomaly detection.
\paragraph{Problem: the definition of {$S_{\text{normal}}$}}
In the original paper a weight vector \(S_\text{normal}\) is required as input and it's described as a representation of normal transaction patterns.
The exact definition was not provided but from the definition of the loss function it's clear that \(S_\text{normal}\) should contain as values the feature vectors \(\phi(G)\) of graphs that are known to be legitimate. This, from my prospective, is a big flaw in the original paper. In fact, the value of \(\phi(G), G \in \text{legitimate}\) cannot be known before the training and therfore it cannot be passed as input to the model. I'm sure that in their implementation they found a way to compute an estimate of \(S_\text{normal}\) and maybe then update it during the traing, but I think that it should have been mentioned in the paper.


To overcome this issue i firstly tried a method inspired by "Deep One-Class Classification" [9] that collapses normal data into a single mean vector (see Appendix \ref{firstmethod}). Unfortunately, the assumption of using a hypersphere to model normal transactions turned out to be too restrictive and the model struggled to generalize to unseen fraud patterns.

By implementing \(S_\text{normal}\) as a Memory Bank we maintain a dynamic collection of feature vectors from the most recent legitimate graphs. This allows the model to represent the ``normal'' class as a multimodal manifold, effectively capturing distinct clusters of valid behavior (e.g., high-frequency small trading vs. low-frequency large transfers) without averaging them into an inaccurate center.

We define the Memory Bank at training step $t$ as a set $\mathcal{M}^{(t)}$ with fixed capacity $K$:
\begin{equation}
    \mathcal{M}^{(t)} \doteq \{ \mathbf{m}_1, \mathbf{m}_2, \dots, \mathbf{m}_K \}
\end{equation}
where each $\mathbf{m}_k \in \mathbb{R}^d$ is a stored feature vector $\phi(G)$ of a previously processed legitimate graph. The bank is updated via a First-In-First-Out (FIFO) queue mechanism. Let $V_{\text{normal}}^{(t)}$ be the set of feature vectors from legitimate graphs in the current batch.

To ensure the unsupervised loss is well-defined at the start of training, the memory bank is initialized using a ``warm-up'' phase. Let $\mathcal{G}_{\text{init}} \subset \mathcal{G}_{\text{normal}}$ be a subset of $K$ legitimate graphs sampled before training begins. We define the initial state $\mathcal{M}^{(0)}$ as:
\begin{equation}
    \mathcal{M}^{(0)} \doteq \{ \phi(G; \Theta^{(0)}) \mid G \in \mathcal{G}_{\text{init}} \}
\end{equation}
where $\Theta^{(0)}$ represents the randomly initialized network parameters. This ensures that the bank contains a representative distribution of the untrained feature space, preventing model collapse during the early training iterations.
Afterwards the update rule is:
\begin{equation}
    \mathcal{M}^{(t+1)} \leftarrow \left( \mathcal{M}^{(t)} \setminus \mathcal{M}_{\text{oldest}} \right) \cup V_{\text{normal}}^{(t)}
\end{equation}
where $\mathcal{M}_{\text{oldest}}$ represents the set of vectors removed to make space for the new entries.

The anomaly detection is formulated as a hypothesis test based on the distance to the \textit{nearest} neighbor in the memory bank, rather than a single center:
\begin{equation}
    \mathcal{H}_{0}: \min_{\mathbf{m} \in \mathcal{M}} ||\phi(G) - \mathbf{m}||^2 \le \tau \quad \text{(Normal)}
\end{equation}
\begin{equation}
    \mathcal{H}_{1}: \min_{\mathbf{m} \in \mathcal{M}} ||\phi(G) - \mathbf{m}||^2 > \tau \quad \text{(Anomaly)}
\end{equation}
The loss function is defined as
\begin{equation}
    \mathcal{L} = \mathcal{L}_{sup}(\phi(G), y; \Theta) + \lambda_{1} \mathcal{L}_{\text{unsup}}(\phi(G), \mathcal{M}) + \lambda_{2}\mathcal{R}(\Theta)
\end{equation}
where the supervised term remains:
\begin{equation}
    \mathcal{L}_{sup}(\phi(G), y; \Theta) = - \left[ y \log(\hat{y}) + (1 - y) \log(1 - \hat{y}) \right]
\end{equation}
with $\hat{y} = \sigma(W^\top \phi(G) + b)$. The unsupervised loss term is redefined to minimize the distance between a normal graph's embedding and its closest prototype in the bank:
\begin{equation}
    \mathcal{L}_{\text{unsup}}(\phi(G), \mathcal{M}) = \min_{\mathbf{m} \in \mathcal{M}} ||\phi(G) - \mathbf{m}||^2
\end{equation}
\begin{equation}
    \mathcal{R}(\Theta)=\sum_{l}||\Theta^{(l)}||_{2}^{2}
\end{equation}
$\lambda_1, \lambda_2$ are hyperparameters controlling the trade-off between supervision, manifold consistency, and regularization. $\Theta$ represent the adaptive parameters optimized via quantum gradient flows.

\begin{equation}
    \Theta^{(t+1)} = \Theta^{(t)} - \eta \nabla_{\Theta} \mathcal{L}\left( \mathcal{N}\left( e^{-i H_{eff}(\Theta)} \rho_{in} e^{i H_{eff}(\Theta)} \right) \right)
\end{equation}
This should lead to a robust learning of normal and fraudulent transaction patterns across quantum and classical domains.

After the training we define \(S_\text{normal}\) using a validation set \(W\) composed by \(M\) non fraudolent graphs as
\begin{equation}
    S_\text{normal} = \{\phi(G^*) \;|\; G^* \in W\}
\end{equation}
\(S_\text{normal}\) will not be modified during the inference phase.

\section{Decision and Interpretability}
The final decision is made based on the distance from the learned normal pattern \(S_\text{normal}\).
\begin{equation}
    s(G) = \min_{\phi(G^*) \in S_\text{normal}} ||\phi(G) - \phi(G^*)||^2_2
\end{equation}
The decision rule is
\begin{equation}
    \text{Fraud}(G) \doteq \begin{cases}
        1 & \text{if } s(G) > \tau \\
        0 & \text{otherwise}
    \end{cases}
\end{equation}
Where \(\tau\) is a threshold that can be tuned based on the desired trade-off between false positives and false negatives.
My implementation is slightly different from the one described in the original paper since I used a different definition for \(S_\text{normal}\).
